% ==============================================================================
% Fichier : chapitres/09_securite_applications.tex
% Livre IX : Sécurité des Applications
% ==============================================================================

\chapter{Sécurité des Applications}
\label{chap:securite_apps}

\begin{boxobjectifs}
\textbf{Objectifs du Livre~IX :} Maîtriser les méthodes de développement sécurisé (SSDLC), l'évaluation de la sécurité des logiciels (boîte noire/grise/blanche), les bases de données (SQL, NoSQL, injections), la programmation orientée objet sécurisée et les principes du code défensif.
\end{boxobjectifs}

% ------------------------------------------------------------------------------
\section{Méthodes de Développement}
% ------------------------------------------------------------------------------

\subsection{Méthodes Classiques}
Les méthodes séquentielles (cascade, modèle en~V, itératif) imposent une planification détaillée et une documentation complète à chaque phase. Elles manquent de flexibilité face aux changements tardifs d'exigences, notamment les exigences de sécurité découvertes en fin de cycle.

\subsection{Méthodes Agiles}
Scrum, XP (Extreme Programming) et Kanban privilégient des itérations courtes (sprints), des livraisons fréquentes et l'adaptation continue. Elles facilitent l'intégration de la sécurité à chaque sprint (\textit{Security by Design}).

\subsection{SSDLC (Secure Software Development Life Cycle)}
\begin{boxdef}
Le \textbf{SSDLC} est un cycle de développement logiciel qui intègre la sécurité à \textit{chaque étape} : de la définition des exigences à la mise en production et à la maintenance. L'objectif est de «~déplacer la sécurité vers la gauche~» (\textit{Shift Left Security}) pour détecter les vulnérabilités le plus tôt possible, là où leur correction coûte le moins cher.
\end{boxdef}

\begin{table}[H]
\centering
\small
\begin{tabular}{p{3.5cm}p{8cm}}
    \toprule
    \textbf{Phase SSDLC} & \textbf{Activité de Sécurité} \\
    \midrule
    Exigences         & Définir les exigences de sécurité (authentication, autorisation, chiffrement). \\
    Conception        & Modélisation des menaces (STRIDE, DREAD), architecture sécurisée. \\
    Développement     & Revue de code sécurisé, respect des bonnes pratiques (OWASP). \\
    Test              & Tests de sécurité (SAST, DAST, pentest, fuzzing). \\
    Déploiement       & Configuration sécurisée, scan de vulnérabilités pré-production. \\
    Maintenance       & Gestion des patchs, surveillance continue, réponse aux incidents. \\
    \bottomrule
\end{tabular}
\caption{Activités de sécurité par phase SSDLC}
\end{table}

% ------------------------------------------------------------------------------
\section{Évaluation de la Sécurité des Logiciels}
% ------------------------------------------------------------------------------

\subsection{Vulnérabilités Logicielles Communes}

\begin{table}[H]
\centering
\small
\begin{tabular}{p{4cm}p{8cm}}
    \toprule
    \textbf{Vulnérabilité} & \textbf{Description} \\
    \midrule
    \textbf{Injection de code} & Exploitation d'une entrée utilisateur pour injecter du code (SQL, shell, JavaScript) dans l'application. \\
    \textbf{Exposition de données sensibles} & Exposition involontaire de mots de passe, clés d'API, informations personnelles. \\
    \textbf{Dépassement de tampon (Buffer Overflow)} & Écriture au-delà de la mémoire allouée permettant l'exécution de code arbitraire. \\
    \textbf{Vulnérabilités des composants tiers} & Bibliothèques/frameworks présentant des CVE connus non patchés. \\
    \textbf{Manque de validation des entrées} & Absence de filtrage des données saisies par l'utilisateur. \\
    \textbf{Mauvaise gestion des erreurs} & Messages d'erreur trop verbeux révélant la technologie ou l'architecture interne. \\
    \bottomrule
\end{tabular}
\caption{Vulnérabilités logicielles courantes}
\end{table}

\subsection{Types de Tests de Sécurité Logicielle}

\begin{itemize}
    \item \textbf{SAST (Static Application Security Testing) :} Analyse du code source sans exécution. Détecte les vulnérabilités tôt dans le cycle.
    \item \textbf{DAST (Dynamic Application Security Testing) :} Tests en boîte noire sur l'application en cours d'exécution (ex. : BurpSuite, OWASP ZAP).
    \item \textbf{IAST (Interactive AST) :} Combinaison de SAST et DAST via des agents instrumentant l'application pendant les tests.
    \item \textbf{Fuzzing :} Envoi de données aléatoires ou malformées pour provoquer des comportements inattendus.
\end{itemize}

\subsection{Méthodologie de Test de Sécurité Applicative}
\begin{enumerate}
    \item \textbf{Reconnaissance :} Identification de la technologie, des endpoints et de l'architecture (Maltego, Shodan, Nmap, Burp Spider).
    \item \textbf{Recherche des vulnérabilités :} OWASP~ZAP, SQLMap, DirBuster, Nikto.
    \item \textbf{Exploitation :} Metasploit, BurpSuite, BeEF (XSS).
    \item \textbf{Élévation de privilèges :} Mimikatz, PowerSploit, Empire.
    \item \textbf{Reporting :} Fiches CVSS, impact business, recommandations de correction.
\end{enumerate}

% ------------------------------------------------------------------------------
\section{Sécurité des Bases de Données}
% ------------------------------------------------------------------------------

\subsection{Bases de Données SQL}
Les SGBD relationnels (MySQL, PostgreSQL, Oracle, SQL Server, MariaDB) organisent les données en tables liées par des clés primaires et étrangères. Les données sont interrogées via le langage SQL.

\begin{boxalert}
\textbf{Injection SQL :} Vulnérabilité la plus critique des applications web (OWASP \#1). Elle survient lorsque des entrées utilisateur non filtrées sont incorporées dans une requête SQL.

\begin{lstlisting}[language=SQL]
-- Requête vulnérable
SELECT * FROM users WHERE username='$input' AND password='$pass';

-- Payload d'attaque : ' OR '1'='1' --
-- La requête devient :
SELECT * FROM users WHERE username='' OR '1'='1' --' AND password='...';
-- Résultat : toujours vraie -> accès accordé sans mot de passe
\end{lstlisting}

\textbf{Contre-mesure :} Toujours utiliser des \textbf{requêtes préparées (prepared statements)} :
\begin{lstlisting}[language=PHP]
$stmt = $pdo->prepare(
    "SELECT * FROM users WHERE username = ? AND password = ?"
);
$stmt->execute([$username, $hashed_password]);
\end{lstlisting}
\end{boxalert}

\subsection{Bases de Données NoSQL}
Les SGBD non relationnels (MongoDB, Redis, Cassandra, CouchBase) utilisent des modèles de données flexibles :

\begin{table}[H]
\centering
\small
\begin{tabular}{p{3cm}p{4cm}p{4.5cm}}
    \toprule
    \textbf{Type NoSQL} & \textbf{Modèle de données} & \textbf{Exemples de SGBD} \\
    \midrule
    Clé-valeur       & Paires clé $\Rightarrow$ valeur simples. & Redis, Riak \\
    Orienté document & Documents JSON / XML. & MongoDB, CouchBase \\
    Orienté colonnes & Familles de colonnes. & Cassandra, HBase \\
    Graphe           & Nœuds et relations. & Neo4j, Amazon Neptune \\
    \bottomrule
\end{tabular}
\caption{Types de bases de données NoSQL}
\end{table}

\subsection{Injections NoSQL}
Les SGBD NoSQL ne sont pas immunisés contre les injections. Exemple MongoDB :
\begin{lstlisting}[language=javascript, caption=Injection NoSQL MongoDB]
// Requête vulnérable (Node.js / Express)
db.users.find({ username: req.body.user, password: req.body.pass });

// Payload : { "$gt": "" }  -> condition toujours vraie
// username: { "$gt": "" }, password: { "$gt": "" }  => accès accordé
\end{lstlisting}

% ------------------------------------------------------------------------------
\section{Programmation Orientée Objet (POO) et Sécurité}
% ------------------------------------------------------------------------------

\begin{table}[H]
\centering
\small
\begin{tabular}{p{3cm}p{8.5cm}}
    \toprule
    \textbf{Principe POO} & \textbf{Intérêt pour la sécurité} \\
    \midrule
    \textbf{Encapsulation}  & Masque les détails d'implémentation. Protège les attributs sensibles (private/protected). \\
    \textbf{Héritage}       & Réutilisation des contrôles de sécurité dans les classes dérivées. Risque : héritage de vulnérabilités. \\
    \textbf{Polymorphisme}  & Permet de substituer des implémentations de chiffrement sans changer l'interface. \\
    \textbf{Abstraction}    & Interfaces et classes abstraites masquant la complexité cryptographique. \\
    \textbf{Couplage faible} & Composants indépendants : une vulnérabilité dans un module n'expose pas les autres. \\
    \textbf{Cohésion forte} & Un module fait une seule chose bien : facilite la revue et le test de sécurité. \\
    \bottomrule
\end{tabular}
\caption{Principes POO appliqués à la sécurité}
\end{table}

% ------------------------------------------------------------------------------
\section{Programmation Sécurisée}
% ------------------------------------------------------------------------------

\begin{boxinfo}
\textbf{Principes du code défensif (OWASP Secure Coding Practices) :}
\begin{itemize}
    \item \textbf{Valider toutes les entrées :} Whitelist (liste des valeurs acceptées), jamais de confiance implicite.
    \item \textbf{Encoder les sorties :} Prévenir les attaques XSS par échappement HTML/JS.
    \item \textbf{Authentification forte :} Stocker les mots de passe avec bcrypt/Argon2 (jamais MD5 ou SHA-1 seul).
    \item \textbf{Gestion des sessions :} Tokens aléatoires et suffisamment longs, expiration, régénération après login.
    \item \textbf{Contrôle d'accès :} Vérifier les autorisations côté serveur, pas seulement côté client.
    \item \textbf{Cryptographie :} Utiliser les bibliothèques standard (OpenSSL, libsodium), ne jamais implémenter ses propres algorithmes.
    \item \textbf{Gestion des erreurs :} Ne jamais exposer de détails techniques dans les messages d'erreur utilisateur.
    \item \textbf{Journalisation :} Enregistrer les événements de sécurité (tentatives de connexion, accès refusés) sans consigner de données sensibles.
    \item \textbf{Sécurité des composants tiers :} Tenir à jour les dépendances (npm audit, pip check, OWASP Dependency-Check).
\end{itemize}
\end{boxinfo}

% ------------------------------------------------------------------------------
\section{Évaluation des Acquis --- Livre~IX}
% ------------------------------------------------------------------------------

\begin{boxexo}
\textbf{QCM :}
\begin{enumerate}
    \item Le SSDLC est un cycle de développement qui : \textbf{(a) Accélère le développement \quad (b) Intègre la sécurité à chaque étape \quad (c) Se concentre sur la performance \quad (d) Élimine les tests}
    \item Une vulnérabilité logicielle est : \textbf{(a) Une fonctionnalité utile \quad (b) Une faiblesse pouvant être exploitée \quad (c) Une erreur corrigée par un patch \quad (d) Un composant tiers}
    \item La meilleure contre-mesure contre les injections SQL est : \textbf{(a) Utiliser NoSQL \quad (b) Désactiver les logs \quad (c) Utiliser des requêtes préparées \quad (d) Chiffrer la base de données}
    \item L'encapsulation en POO contribue à la sécurité en : \textbf{(a) Accélérant l'exécution \quad (b) Masquant les attributs et implémentations sensibles \quad (c) Permettant l'héritage multiple \quad (d) Simplifiant l'interface}
    \item Le DAST (Dynamic Application Security Testing) teste : \textbf{(a) Le code source statique \quad (b) L'application en cours d'exécution \quad (c) La base de données uniquement \quad (d) L'infrastructure réseau}
    \item Le principe «~Shift Left Security~» signifie : \textbf{(a) Déplacer les tests en fin de projet \quad (b) Intégrer la sécurité dès les premières phases du développement \quad (c) Sécuriser uniquement l'interface utilisateur \quad (d) Réduire le budget sécurité}
\end{enumerate}

\textbf{Travail Pratique :}\\
Réaliser une évaluation complète de la sécurité de l'application de gestion des notes de l'ENSPY :
\begin{enumerate}
    \item Identifier les vulnérabilités OWASP~Top~10 présentes.
    \item Tester l'injection SQL sur les formulaires de connexion et de recherche.
    \item Proposer et implémenter les corrections (requêtes préparées, validation des entrées, hachage bcrypt).
    \item Rédiger un rapport avec fiches CVSS pour chaque vulnérabilité identifiée.
\end{enumerate}
\end{boxexo}
